\documentclass[11pt,a4paper]{article}

\usepackage[utf8]{inputenc}
\usepackage[T1]{fontenc}
\usepackage{amsmath,amssymb}
\usepackage{graphicx}
\usepackage[margin=1in]{geometry}
\usepackage{hyperref}
\usepackage{booktabs}
\usepackage{caption}
\usepackage{subcaption}
\usepackage[dvipsnames]{xcolor}

% Path to figures
\graphicspath{{../phase2d_r5_e50_p0.9000-0.9950_b16-60000_n5000-60000_w512_lr0.01_logtr100_logte500/figures_notebook/}}

\title{Phase Diagram of Noisy MNIST Classification:\\
  A Finite-Size Scaling Analysis of Corruption, Batch Size, and Dataset Size}
\author{}
\date{}

\begin{document}
\maketitle

%% ============================================================
\section{Introduction}
%% ============================================================

A fundamental question in supervised learning is: \emph{under what conditions can a neural network extract useful signal from heavily corrupted training data?}
This report documents a systematic computational study that maps out the \textbf{phase diagram} of a single-hidden-layer MLP trained on MNIST digits whose pixels have been subjected to salt-and-pepper noise.
The analysis draws on concepts from statistical physics---phase transitions, critical exponents, and finite-size scaling---to characterize the sharp boundary between a ``learning'' phase (where the network achieves non-trivial test accuracy on clean data) and a ``non-learning'' phase (where performance remains at chance level, ${\sim}10\%$).

The experiments sweep over three control parameters:
\begin{enumerate}
  \item \textbf{Corruption probability} $p \in [0.90,\, 0.995]$: the fraction of input pixels replaced by uniform random values.
  \item \textbf{Batch size} $B \in \{16, 32, 64, \ldots, 4096, 60000\}$: spanning minibatch SGD through full-batch gradient descent.
  \item \textbf{Training set size} $N \in \{5000,\, 10000,\, 20000,\, 40000,\, 60000\}$: used to probe finite-size scaling behavior.
\end{enumerate}
For each $(p, B, N)$ triplet the experiment is repeated over 5 random seeds, and the mean clean-test accuracy is recorded as the order parameter.

%% ============================================================
\section{Experimental Setup}
%% ============================================================

\subsection{Model Architecture}

The classifier is a single-hidden-layer MLP:
\[
  f(x) = W_2 \,\sigma\!\bigl(W_1 \,\mathrm{vec}(x) + b_1\bigr) + b_2,
\]
where $x \in \mathbb{R}^{28\times28}$ is the input image, $W_1 \in \mathbb{R}^{h \times 784}$, $W_2 \in \mathbb{R}^{10 \times h}$, $\sigma$ is the ReLU activation, and the hidden width is $h = 512$.

\subsection{Noise Model}

Each training image undergoes \textbf{salt-and-pepper corruption}: every pixel is independently replaced with a sample from $\mathrm{Uniform}(0,1)$ with probability $p$, and left unchanged with probability $1 - p$.
At $p = 0.90$, $90\%$ of pixels are noise; at $p = 0.995$, only $0.5\%$ of the original signal survives.
Crucially, the test set is always evaluated \emph{clean} (no corruption), so test accuracy measures whether the network has extracted genuine digit structure from the noisy training signal.

\subsection{Training Protocol}

All runs use SGD with learning rate $\eta = 0.01$ and no momentum or weight decay.
The total computation budget is fixed at $N \times 50$ sample presentations (i.e., 50 ``reference epochs''), regardless of batch size.
This means that a run with $B = 64$ takes many more gradient steps than one with $B = 60000$ (full-batch), ensuring a fair comparison of the \emph{information content per sample} rather than per step.

\subsection{Parameter Grid}

\begin{table}[h]
\centering
\begin{tabular}{ll}
\toprule
\textbf{Parameter} & \textbf{Values} \\
\midrule
Corruption $p$ & 30 values linearly spaced in $[0.90, 0.995]$ \\
Batch size $B$ & $16, 32, 64, 128, 256, 512, 1024, 2048, 4096, 60000$ \\
Dataset size $N$ & $5000,\; 10000,\; 20000,\; 40000,\; 60000$ \\
Seeds & $0, 1, 2, 3, 4$ \\
\bottomrule
\end{tabular}
\caption{Experiment parameter grid. Total configurations: $30 \times 10 \times 5 \times 5 = 7{,}500$ individual runs.}
\label{tab:params}
\end{table}

%% ============================================================
\section{Results and Observations}
%% ============================================================

\subsection{Phase Diagram at Full Dataset Size}

Figure~\ref{fig:heatmap_N60000} shows the primary result: a heatmap of mean clean-test accuracy in the $(p, B)$ plane at $N = 60000$.

\begin{figure}[ht]
  \centering
  \includegraphics[width=0.75\textwidth]{heatmap_N60000_notebook.png}
  \caption{Phase diagram at $N = 60000$. Color encodes mean clean-test accuracy (blue $= 1.0$, red $= 0.1$ chance level). The black contour marks the $15\%$ accuracy threshold. A sharp boundary separates the learning phase (upper-left) from the non-learning phase (lower-right).}
  \label{fig:heatmap_N60000}
\end{figure}

\textbf{Key observations:}
\begin{itemize}
  \item For small batch sizes ($B \leq 512$), the network can learn effectively even when ${\sim}98\%$ of pixels are corrupted. The critical corruption probability $p_c$ remains close to $0.995$.
  \item As batch size increases beyond ${\sim}1024$, the learning region shrinks sharply. At full-batch ($B = 60000$), learning fails at much lower corruption levels ($p_c \approx 0.98$).
  \item The transition is not gradual---accuracy drops rapidly from ${\sim}70\%$ to chance level over a narrow band of $p$ values, suggesting a genuine phase transition rather than a smooth degradation.
\end{itemize}

This reveals a striking result: \textbf{minibatch SGD noise acts as an implicit regularizer that extends learnability into regimes where full-batch training fails.}

\subsection{Effect of Dataset Size}

Figure~\ref{fig:heatmap_by_N} shows the phase diagram at each dataset size $N$.

\begin{figure}[ht]
  \centering
  \includegraphics[width=\textwidth]{heatmap_by_N_notebook.png}
  \caption{Phase diagrams for $N \in \{5000, 10000, 20000, 40000, 60000\}$. As $N$ increases, the learning region expands (the phase boundary shifts rightward to higher $p$ and downward to larger $B$), and the transition becomes sharper.}
  \label{fig:heatmap_by_N}
\end{figure}

\textbf{Key observations:}
\begin{itemize}
  \item At $N = 5000$, only small batch sizes at moderate corruption can learn. The learning region is narrow.
  \item With each doubling of $N$, the phase boundary shifts noticeably, indicating that more data helps the network tolerate greater corruption.
  \item The transition becomes progressively sharper with increasing $N$---consistent with finite-size scaling theory, where the system size (here, dataset size) controls the sharpness of the transition.
\end{itemize}

\subsection{Full-Batch vs.\ Minibatch Comparison}

Figure~\ref{fig:fullbatch_vs_minibatch} directly compares accuracy--vs--$p$ curves for several batch sizes at $N = 60000$.

\begin{figure}[ht]
  \centering
  \includegraphics[width=0.7\textwidth]{full_batch_vs_minibatch_notebook.png}
  \caption{Mean clean-test accuracy as a function of corruption probability $p$ for batch sizes $B \in \{64, 256, 1024, 60000\}$ at $N = 60000$. Smaller batches sustain higher accuracy deeper into the high-corruption regime.}
  \label{fig:fullbatch_vs_minibatch}
\end{figure}

\textbf{Key observations:}
\begin{itemize}
  \item At $B = 64$, accuracy remains above $70\%$ up to $p \approx 0.97$ and above $50\%$ up to $p \approx 0.98$.
  \item Full-batch training ($B = 60000$) already begins degrading at $p \approx 0.93$ and collapses to chance by $p \approx 0.99$.
  \item There is a clear ordering: smaller $B$ $\Rightarrow$ higher noise tolerance. The gap between curves widens as $p$ increases, indicating that the regularization benefit of SGD noise grows in the high-corruption regime.
\end{itemize}

\subsection{Phase Boundary and Power-Law Fit}

Figure~\ref{fig:phase_boundary} plots the estimated critical corruption $p_c(B)$ at $N = 60000$ for two accuracy thresholds.

\begin{figure}[ht]
  \centering
  \includegraphics[width=0.7\textwidth]{phase_boundary_notebook.png}
  \caption{Estimated critical corruption probability $p_c$ as a function of batch size at $N = 60000$, for accuracy thresholds $0.15$ and $0.20$. Dashed lines show power-law fits of the form $p_c(B) = 1 - a \, B^{-\beta}$.}
  \label{fig:phase_boundary}
\end{figure}

The phase boundary is fit to the functional form
\[
  p_c(B) \;=\; 1 - a\, B^{-\beta},
\]
or equivalently, the critical noise fraction scales as $(1 - p_c) \propto B^{-\beta}$.
This power-law relationship suggests that the critical corruption tolerable by the network decreases algebraically with batch size.

\textbf{Key observations:}
\begin{itemize}
  \item For small batch sizes ($B \lesssim 256$), $p_c$ is nearly flat near $0.995$---the network can learn even with $99.5\%$ corruption.
  \item Above $B \approx 1000$, $p_c$ drops sharply, indicating a regime change where SGD noise is no longer sufficient to regularize against input corruption.
  \item The power-law fit captures the overall trend, though the sharp plateau at small $B$ suggests the true functional form may involve a saturation mechanism.
\end{itemize}

\subsection{Training Dynamics Near the Phase Boundary}

Figure~\ref{fig:training_dynamics} shows training loss and clean-test accuracy trajectories for two representative runs near the phase boundary.

\begin{figure}[ht]
  \centering
  \includegraphics[width=0.7\textwidth]{training_dynamics_notebook.png}
  \caption{Training dynamics for two runs near the phase boundary ($N = 60000$). Top: training loss. Bottom: clean-test accuracy. Both runs hover near chance-level loss ($\ln 10 \approx 2.303$) and accuracy (${\sim}10$--$17\%$), exhibiting the marginal, noisy learning characteristic of the critical region.}
  \label{fig:training_dynamics}
\end{figure}

\textbf{Key observations:}
\begin{itemize}
  \item Training loss remains very close to $\ln(10) \approx 2.303$, the cross-entropy of a uniform predictor over 10 classes. This is the hallmark of near-chance performance.
  \item Test accuracy fluctuates between ${\sim}10\%$ and ${\sim}17\%$, only marginally above chance. The network extracts a faint signal but cannot robustly learn.
  \item The curves are noisy and non-monotone, reflecting the stochastic competition between signal and noise at the critical point.
\end{itemize}

\subsection{Finite-Size Scaling}

Figure~\ref{fig:fss_curves} shows accuracy vs.\ $p$ curves at fixed batch sizes for varying $N$, the analog of varying system size in a statistical-physics finite-size scaling analysis.

\begin{figure}[ht]
  \centering
  \includegraphics[width=\textwidth]{finite_size_scaling_curves_notebook.png}
  \caption{Finite-size scaling curves for $B \in \{64, 256, 1024\}$. Each panel shows mean clean-test accuracy vs.\ $p$ for $N \in \{5000, 10000, 20000, 40000, 60000\}$. The curves steepen and converge toward a common crossing point as $N$ increases.}
  \label{fig:fss_curves}
\end{figure}

\textbf{Key observations:}
\begin{itemize}
  \item For each batch size, the accuracy curves for different $N$ values approximately cross at a common critical point $p_c$. This crossing behavior is a hallmark of a continuous phase transition.
  \item At $B = 64$, the crossing point is near $p_c \approx 0.985$. At $B = 1024$, it shifts to $p_c \approx 0.975$.
  \item The transition sharpens with increasing $N$: at $N = 60000$ the drop from high accuracy to chance is nearly vertical, while at $N = 5000$ it is a broad crossover.
  \item Larger $N$ consistently yields higher accuracy at any fixed $p$, confirming that more data helps overcome corruption.
\end{itemize}

\subsection{Finite-Size Scaling Collapse}

Figure~\ref{fig:fss_collapse} tests whether the data from different $N$ values collapse onto a single universal curve under the finite-size scaling ansatz:
\[
  A(p, N) \;=\; A_c + N^{-\kappa}\, \mathcal{F}\!\bigl((p - p_c)\, N^{1/\nu}\bigr),
\]
where $A$ is accuracy, $p_c$ is the critical corruption probability, $\nu$ is the correlation-length exponent, $\kappa$ governs the scaling of the order parameter at criticality, $A_c$ is the critical accuracy, and $\mathcal{F}$ is a universal scaling function.

\begin{figure}[ht]
  \centering
  \includegraphics[width=\textwidth]{finite_size_collapse_notebook.png}
  \caption{Finite-size scaling collapse for $B \in \{64, 256, 1024\}$. Rescaled axes: $x = (p - p_c)\,N^{1/\nu}$, $y = (A - A_c)\,N^{\kappa}$. The fitted critical parameters $(p_c, \nu, A_c, \kappa)$ are shown in each panel title.}
  \label{fig:fss_collapse}
\end{figure}

\textbf{Key observations:}
\begin{itemize}
  \item The collapse is best at $B = 256$ and $B = 1024$, where curves from all five $N$ values approximately overlap onto a single sigmoid-like scaling function.
  \item The fitted critical corruption probability $p_c$ ranges from ${\sim}0.986$ ($B = 64$) to ${\sim}0.989$ ($B = 1024$), consistent with the crossing-point estimates from the raw curves.
  \item The correlation-length exponent $\nu$ ranges from ${\sim}0.56$ to ${\sim}0.80$ across the three batch sizes. A smaller $\nu$ indicates a sharper transition in the scaling variable.
  \item The exponent $\kappa = 2.0$ was found at the boundary of the search grid for all batch sizes, suggesting that the accuracy at criticality decays as $N^{-2}$ and that the true optimal $\kappa$ may lie outside the search range.
  \item The critical accuracy $A_c$ ranges from ${\sim}0.20$ to ${\sim}0.56$, reflecting the batch-size dependence of the order parameter at the transition.
\end{itemize}

%% ============================================================
\section{Technical Description of the Experiment Script}
%% ============================================================

The experiment is implemented in a single self-contained Python script (\texttt{run\_experiments\_2d\_phase\_diagram.py}) with the following structure.

\subsection{Data Pipeline}

\begin{enumerate}
  \item \textbf{Base dataset}: MNIST is loaded via \texttt{torchvision} with a simple \texttt{ToTensor()} transform.
  \item \textbf{Subset selection}: For $N < 60000$, a \emph{balanced} subset is drawn: exactly $N/10$ examples per digit class, using a deterministic seed for reproducibility.
  \item \textbf{Corruption}: A custom \texttt{SaltAndPepperDataset} wrapper applies per-pixel corruption on-the-fly. Each pixel is replaced with $\mathrm{Uniform}(0,1)$ with probability $p$, using a seeded generator for reproducibility.
  \item \textbf{Evaluation}: The clean MNIST test set (no corruption) is used for all accuracy measurements.
\end{enumerate}

\subsection{Training Loop}

The training budget is defined in terms of \emph{total samples seen}:
\[
  T = N \times E_{\mathrm{ref}} = N \times 50.
\]
The number of gradient steps is then $\lceil T / B_{\mathrm{eff}} \rceil$, where $B_{\mathrm{eff}} = \min(B, N)$ is the effective batch size (capped at the dataset size).
This ensures that changing $B$ does not change the total amount of data the network sees---only the granularity of updates.

The optimizer is vanilla SGD with fixed learning rate $\eta = 0.01$, no momentum, no weight decay, and no learning rate schedule.

\subsection{Parallelism and Execution}

The script uses Python's \texttt{ProcessPoolExecutor} to run experiments in parallel across CPU cores, with configurable worker counts and CPU affinity pinning. Each worker caches the MNIST dataset in a process-global dictionary to avoid redundant I/O. Results for each $(p, B, N, \text{seed})$ combination are atomically written to individual JSON files, enabling crash-safe resumption.

\subsection{Aggregation and Plotting}

After all runs complete, results are aggregated by computing mean and standard error of test accuracy across seeds for each $(p, B, N)$ triplet. The script then generates the following diagnostic plots:
\begin{itemize}
  \item \textbf{Heatmaps}: 2D color maps of accuracy in the $(p, B)$ plane, with contour lines at the phase threshold.
  \item \textbf{Phase boundary}: Estimated $p_c(B)$ via linear interpolation of accuracy crossing a threshold, with power-law fits.
  \item \textbf{Finite-size scaling curves}: Accuracy vs.\ $p$ at fixed $B$ for varying $N$.
  \item \textbf{Scaling collapse}: Rescaled data using optimized critical exponents $(p_c, \nu, \kappa)$, found via grid search minimizing the bin-variance collapse score.
  \item \textbf{Training dynamics}: Loss and accuracy trajectories for representative runs near the phase boundary.
  \item \textbf{Full-batch comparison}: Direct comparison of minibatch vs.\ full-batch accuracy curves.
\end{itemize}

\subsection{Scaling Collapse Methodology}

Two collapse methods are implemented:

\paragraph{Custom grid search.}
The script performs a brute-force search over $(p_c, \nu, \kappa)$ to minimize a \emph{collapse quality score}: the data is binned along the scaled $x$-axis and the mean within-bin variance of the scaled $y$-values is computed. Lower variance $\Rightarrow$ better collapse. The search grid is:
\begin{itemize}
  \item $p_c$: 121 values in a narrow window around the initial estimate.
  \item $\nu$: 49 values in $[0.2, 5.0]$ (or a refined range around the initial estimate).
  \item $\kappa$: 81 values in $[-2.0, 2.0]$.
\end{itemize}

\paragraph{pyfssa library.}
An alternative collapse is attempted using the \texttt{pyfssa} library (finite-size scaling analysis), which implements a more principled quality-of-collapse metric. This method fixes the anomalous dimension $\zeta = 0$ and searches over $(p_c, \nu)$ only. In the current experiment run, this method failed due to a NumPy compatibility issue.

%% ============================================================
\section{Summary}
%% ============================================================

This experiment establishes that the transition from learning to non-learning under heavy pixel corruption in MNIST exhibits the hallmarks of a \emph{continuous phase transition}:
\begin{itemize}
  \item A sharp boundary in the $(p, B)$ control-parameter space.
  \item Finite-size scaling: the transition sharpens as dataset size $N$ increases.
  \item Approximate data collapse under a scaling ansatz with well-defined critical exponents.
  \item A common crossing point in accuracy curves at different $N$, identifying the critical corruption $p_c$.
\end{itemize}

The most striking practical finding is the role of batch size: \textbf{smaller batches dramatically extend the learnable regime}, allowing networks to extract signal even when $99\%+$ of pixels are noise.
This suggests that the implicit regularization provided by minibatch SGD stochasticity is not merely a convenience but a qualitatively important factor that shifts the phase boundary of learnability.

\end{document}
